\chapter[Sermon 56. The Beast Is Worshipped, and He Blasphemes the Name of God, and the Saints of God, and Finally Makes War with the Saints.]{The Beast Is Worshipped, and He Blasphemes the Name of God, and the Saints of God, and Finally Makes War with the Saints.}
\label{sermon:056}
\markright{Sermon 56. The Beast Is Worshipped, and He Blasphemes the Name of ...}

And they worshipped the beast, saying, “Who is like unto the beast? Who is able to war with him?” And there was given to him a mouth to speak great things, and blasphemies: and power was given unto him, to do forty-two months. And he opened his mouth unto blasphemy against God, to blaspheme his name, and his tabernacle, and them that dwell in heaven. And it was given unto him to make war with the saints and to overcome them.

\index{John, Saint}\index{Papacy}\index{Rome}He states that the world worships the dragon; now he adds that it worships the beast. If it is said that the beast represents the empire, some might wonder how the empire could be worshipped. But we will explain briefly how people worship the empire: by accepting its decrees, rituals, and superstitious ordinances, and by being entirely dependent on it. And there were not a few at that time who, in order to gain favor with the Roman Empire, denied the faith of Christ, revolted from the church, and joined themselves in religious practices and sacrifices to the fellowship of the empire. They truly worshipped the beast. Furthermore, that which is due only to one God, the Romans attributed to their empire. But whoever ascribes divine properties to anything, does in truth deify and worship it. And the properties of God are these: to be without equal or end, that he alone is greatest and best, immortal, eternal, most powerful, most invincible. For so the Prophets say: Who is like unto you, O God, in heaven and in earth? Who is as you? Who can resist God? But the Romans attributed all these things to their emperors and to their empire, saying, as Saint John also recounts: Who is like unto Rome? Who is able to war with it? They called their emperors gods, best, greatest, most powerful, and most invincible. They called the empire itself eternal. You can still see these things in ancient authors and coins. Therefore, those who were not ashamed to attribute these things to the Roman rulers and kingdom rightly worshipped the beast. And what other thing is done today, while for the favor of emperors, kings, popes, and their realms, the truth is denied or twisted according to the affections of people? These also worship the beast.

\index{Daniel (Book of)}Now the beast is also given a mouth, speaking great things and blasphemies. Of blasphemies we shall speak more fully later. But since the Roman Empire obtained the greatest victories and held the most magnificent triumphs, it seems to have been given opportunity to boast proudly of those victories, and to claim as their own things that were in truth accomplished through the power of God. And without doubt, the greatest and most unrestrained boasts of the Romans remain, that they are conquerors and lords of the world. But such pride was severely punished in Nebuchadnezzar the king, as you may see in chapter 4 of Daniel. Saint Peter affirms that God resists the proud and gives grace to the humble. God hates the arrogant and takes away their names from the earth.

And where some person might demand, what will be the end of injuries, pride, and finally, intolerable authority and blasphemies? John anticipates this and says, “And power was given him to do, that is to work violence, for forty-two months,” that is to say, for such a time as seems good to the Lord. Although he wishes the time to be unknown to us, it is known to him, so that the faithful may promise themselves that this evil will endure only a few months. I have reasoned about this number in Chapter 11 and Sermon 46, and have shown in previous passages that those numbers are equivalent: namely, the thousand two hundred and three score days, the forty-two months, the time, two times, and half a time. God, therefore, admonishing us as it were by a riddle, does not want us to curiously inquire after times which he has kept in his own power; it is sufficient for us that he has assigned all things within their proper limits.

\index{Augustine, Saint}A plentiful treatise of Roman blasphemies. First, he says by a trope, he has opened his mouth: whereby he has signified his boldness, and liberty, yea licentiousness of speaking. For we say he would not once open his mouth, when we signify any man that will not speak frankly. But the Romans, and companions of the Roman superstition, blaspheme God in three manners. For first, they blaspheme the holy name of God in this, that they do prefer their false gods and their superstitions to the true God, to the true and most holy religion. For where they did admit in the city of Rome the gods of all nations, and their religions, the religion of the only God of Israel they utterly refused: for that they understood how he would be worshipped alone, and by no other rite than that which he himself had prescribed. But they had rather retain wickedly those their many gods, and their religion, although most absurd, than to commit themselves into the tuition of one, and to receive a moderate and simple religion. Augustine, Aurel. I recount not now the blasphemous words of them, uttered against the true God, about that time chiefly, when Vespasian and Titus triumphed, after the Jewish war finished, both of the City destroyed, and the people of God overcome. There were carried about in the triumph the holy vessels of the Temple, and even the God of the Jews, as vanquished and bounden, was seen led into the Capitol house, to make his supplication to their great god Jupiter, as it pleased them. Whereupon we understand that the name of God was no whit less outrageously blasphemed at that time than it was in old time of the Palestinians or Philistines, what time they set the Ark in the temple of their god Dagon; likewise of Rehoboam [?] and Sennacherib, moreover of Belshazzar, King of Babylon, in the 5th chapter of Daniel. But the offenders are found out at the last.

Secondly, the Romans blasphemed the Tabernacle of God. That same ancient Tabernacle of the people of Israel was not only the office, or place of religion and worship, but also a token of God’s presence. For God is now present in the midst of his Church, a figure of whom the Tabernacle of witness represented. But the Romans called the Christian church wicked, foolish, seditious, harlike, and detestable; and they most grievously persecuted it, seeking to destroy it by all means, and bent their whole power to this end.

Finally, they also blasphemed the heavenly inhabitants, the blessed souls of saints, prophets, and apostles, whom they called wicked, seducers, peace breakers, blasphemers, heretics, and sinful persons. For at this time, while John wrote these things, several apostles under the Roman Empire had already been executed and slain, considered plagues of the world, and their memory and doctrine condemned. But from this, you perceive how displeased God is when anyone reviles godly preachers and holy ministers of the church. For the Lord takes the reproach spoken as if against himself. There remain even today blasphemies of this sort, as found in Cornelius Tacitus in the twenty-first book of Augustine, written against Moses and the people of God.

\index{Jerome, Saint}\index{Eusebius of Caesarea}Moreover, God permits the beast to wage war against the saints and to overcome them. For the Roman Empire, until the time of Constantine the Great, instigated ten most severe persecutions against the church. You may read of this in Eusebius, bishop of Caesarea, and Orosius in the history he wrote to Jerome. And this passage especially pertains to the instruction and comfort of the church. For the Lord also in the Gospel prophesies of the church’s destinies, for the consolation and information of the faithful, as appears in chapters 15 and 16 of John. And how the saints are overcome I explained in chapter 11. The Lord Jesus preserve his church. Amen.
